\section{Arm CCA I/O、DMA、accelerator、interrupt}

\begin{frame}[t,shrink=6]{Arm CCA I/O、DMA、accelerator、interrupt：方向开场}
\DeckKicker{方向开场}
\SlideMeta{证据等级}{方向综述}{来源状态：公开材料综合}
{\large\bfseries\color{Ink} Arm CCA I/O、DMA、accelerator、interrupt 的核心是先界定保护对象、管理边界和证据强度，再用三篇主讲材料串起技术演进。\par}\vspace{1.2mm}
\begin{columns}[T,totalwidth=\textwidth]
  \begin{column}{0.45\textwidth}
    \fcolorbox{Rule}{Paper}{%
  \begin{minipage}[t]{0.97\linewidth}
    \vspace{1.1mm}
    \hspace{1.4mm}\begin{minipage}{0.92\linewidth}
      \PanelTitle{内容说明}
      \scriptsize \begin{itemize}
  \item 把设备访问、DMA、interrupt ownership 与 accelerator workflowow 纳入 CCA 边界。
  \item 固定三篇主讲材料；ACAI 和 Devlore 均保持 draft/not-ratified 状态，`sok-tee` 只作为 accelerator TEE taxonomy。
  \item 先讲方向痛点和设计轴，再逐篇说明贡献、限制、证据强度和可落地条件。
  \item 对规范、Survey、draft、vendor evidence 保持显式标注，避免把背景材料写成一手系统证明。
\end{itemize}
    \end{minipage}
    \vspace{1mm}
  \end{minipage}}
  \end{column}
  \begin{column}{0.49\textwidth}
    \fcolorbox{Rule}{Panel}{%
  \begin{minipage}[t]{0.97\linewidth}
    \vspace{1.1mm}
    \hspace{1.4mm}\begin{minipage}{0.92\linewidth}
      \PanelTitle{三篇主讲材料的分工}
      \scriptsize {\tiny
\begin{tabularx}{\linewidth}{@{}YYY@{}}
\toprule
\textbf{论文} & \textbf{定位} & \textbf{证据边界} \\
\midrule
ACAI: CCA accelerator execution & 开创/基础入口 & E3 / 草案/未批准规范 \\
Devlore: device interrupt protection & 代表性改进 & E3 / 草案/未批准规范 \\
Accelerator TEE SoK for CCA I/O & 当前 SOTA/补充边界 & E2 / Background substrate \\
\bottomrule
\end{tabularx}
}
    \end{minipage}
    \vspace{1mm}
  \end{minipage}}
  \end{column}
\end{columns}
\SourceFootnote{来源：论文原文、官方规范或公开材料｜证据等级：见本页 badge}
\end{frame}


\begin{frame}[t,shrink=6]{ACAI: CCA accelerator execution：内容摘要}
\DeckKicker{内容摘要}
\SlideMeta{证据等级}{E3 / 草案/未批准规范}{来源状态：本地 PDF 已核验}
{\large\bfseries\color{Ink} 研究如何用 Arm CCA 保护 accelerator execution，说明 CPU Realm 之外的 offload path 不能默认。\par}\vspace{1.2mm}
\begin{columns}[T,totalwidth=\textwidth]
  \begin{column}{0.45\textwidth}
    \fcolorbox{Rule}{Paper}{%
  \begin{minipage}[t]{0.97\linewidth}
    \vspace{1.1mm}
    \hspace{1.4mm}\begin{minipage}{0.92\linewidth}
      \PanelTitle{内容说明}
      \scriptsize \begin{itemize}
  \item 研究如何用 Arm CCA 保护 accelerator execution，说明 CPU Realm 之外的 offload path 不能默认继承机密性
  \item ACAI is the early design point for extending Arm CCA toward accelerator execution paths.
\end{itemize}
    \end{minipage}
    \vspace{1mm}
  \end{minipage}}
  \end{column}
  \begin{column}{0.49\textwidth}
    \fcolorbox{Rule}{Panel}{%
  \begin{minipage}[t]{0.97\linewidth}
    \vspace{1.1mm}
    \hspace{1.4mm}\begin{minipage}{0.92\linewidth}
      \PanelTitle{证据定位}
      \scriptsize {\tiny
\begin{tabularx}{\linewidth}{@{}YY@{}}
\toprule
\textbf{字段} & \textbf{内容} \\
\midrule
论文定位 & 开创/基础入口 \\
材料类型 & 系统论文 \\
证据等级 & E3 / 草案/未批准规范 \\
来源状态 & 本地 PDF 已核验 \\
\bottomrule
\end{tabularx}
}
    \end{minipage}
    \vspace{1mm}
  \end{minipage}}
  \end{column}
\end{columns}
\SourceFootnote{来源：ACAI: CCA accelerator execution｜证据等级：E3 / 草案/未批准规范}
\end{frame}


\begin{frame}[t,shrink=6]{ACAI: CCA accelerator execution：研究背景}
\DeckKicker{研究背景}
\SlideMeta{证据等级}{E3 / 草案/未批准规范}{来源状态：本地 PDF 已核验}
{\large\bfseries\color{Ink} Confidential workloads 需要使用 GPU/NPU/FPGA 等加速器\par}\vspace{1.2mm}
\begin{columns}[T,totalwidth=\textwidth]
  \begin{column}{0.45\textwidth}
    \fcolorbox{Rule}{Paper}{%
  \begin{minipage}[t]{0.97\linewidth}
    \vspace{1.1mm}
    \hspace{1.4mm}\begin{minipage}{0.92\linewidth}
      \PanelTitle{内容说明}
      \scriptsize \begin{itemize}
  \item Confidential workloads 需要使用 GPU/NPU/FPGA 等加速器
  \item 如果设备 DMA、队列和驱动仍由 untrusted host 控制，Realm 内数据会在 offload 路径暴露
\end{itemize}
    \end{minipage}
    \vspace{1mm}
  \end{minipage}}
  \end{column}
  \begin{column}{0.49\textwidth}
    \fcolorbox{Rule}{Panel}{%
  \begin{minipage}[t]{0.97\linewidth}
    \vspace{1.1mm}
    \hspace{1.4mm}\begin{minipage}{0.92\linewidth}
      \PanelTitle{问题边界}
      \scriptsize {\tiny
\begin{tabularx}{\linewidth}{@{}YY@{}}
\toprule
\textbf{维度} & \textbf{说明} \\
\midrule
保护对象 & Arm CCA I/O、DMA、accelerator、interrupt \\
传统瓶颈 & 把设备访问、DMA、interrupt ownership 与 accelerator workflowow 纳入 CCA 边界。 \\
本文入口 & ACAI: CCA accelerator execution \\
证据限制 & E3 / 草案/未批准规范 \\
\bottomrule
\end{tabularx}
}
    \end{minipage}
    \vspace{1mm}
  \end{minipage}}
  \end{column}
\end{columns}
\SourceFootnote{来源：ACAI: CCA accelerator execution｜证据等级：E3 / 草案/未批准规范}
\end{frame}


\begin{frame}[t,shrink=6]{ACAI: CCA accelerator execution：解决方案}
\DeckKicker{解决方案}
\SlideMeta{证据等级}{E3 / 草案/未批准规范}{来源状态：本地 PDF 已核验}
{\large\bfseries\color{Ink} 以 CCA-aware accelerator execution path 连接 Realm、设备内存访问、队列管理和 attestation，缩小 host 可见面\par}\vspace{1.2mm}
\begin{columns}[T,totalwidth=\textwidth]
  \begin{column}{0.45\textwidth}
    \fcolorbox{Rule}{Paper}{%
  \begin{minipage}[t]{0.97\linewidth}
    \vspace{1.1mm}
    \hspace{1.4mm}\begin{minipage}{0.92\linewidth}
      \PanelTitle{内容说明}
      \scriptsize \begin{itemize}
  \item 以 CCA-aware accelerator execution path 连接 Realm、设备内存访问、队列管理和 attestation，缩小 host 可见面
\end{itemize}
    \end{minipage}
    \vspace{1mm}
  \end{minipage}}
  \end{column}
  \begin{column}{0.49\textwidth}
    \fcolorbox{Rule}{Panel}{%
  \begin{minipage}[t]{0.97\linewidth}
    \vspace{1.1mm}
    \hspace{1.4mm}\begin{minipage}{0.92\linewidth}
      \PanelTitle{核心设计拆解}
      \scriptsize \begin{itemize}
  \item 01. CCA-aware accelerator execution path
  \item 02. Protected DMA and command/data movement
  \item 03. Driver/runtime TCB boundary for offload
\end{itemize}
    \end{minipage}
    \vspace{1mm}
  \end{minipage}}
  \end{column}
\end{columns}
\SourceFootnote{来源：ACAI: CCA accelerator execution｜证据等级：E3 / 草案/未批准规范}
\end{frame}


\begin{frame}[t,shrink=6]{ACAI: CCA accelerator execution：实验结果}
\DeckKicker{实验结果}
\SlideMeta{证据等级}{E3 / 草案/未批准规范}{来源状态：本地 PDF 已核验}
{\large\bfseries\color{Ink} arXiv 系统论文提供原型/设计评估\par}\vspace{1.2mm}
\begin{columns}[T,totalwidth=\textwidth]
  \begin{column}{0.45\textwidth}
    \fcolorbox{Rule}{Paper}{%
  \begin{minipage}[t]{0.97\linewidth}
    \vspace{1.1mm}
    \hspace{1.4mm}\begin{minipage}{0.92\linewidth}
      \PanelTitle{内容说明}
      \scriptsize \begin{itemize}
  \item arXiv 系统论文提供原型/设计评估
  \item draft/not-ratified 状态下只引用本地 PDF 中已核验的机制方向与结果，不写成标准或生产事实
\end{itemize}
    \end{minipage}
    \vspace{1mm}
  \end{minipage}}
  \end{column}
  \begin{column}{0.49\textwidth}
    \fcolorbox{Rule}{Panel}{%
  \begin{minipage}[t]{0.97\linewidth}
    \vspace{1.1mm}
    \hspace{1.4mm}\begin{minipage}{0.92\linewidth}
      \PanelTitle{实验/证据边界}
      \scriptsize {\tiny
\begin{tabularx}{\linewidth}{@{}YYY@{}}
\toprule
\textbf{对象} & \textbf{可支撑} & \textbf{边界} \\
\midrule
数据/来源 & ACAI: CCA accelerator execution & 本地 PDF 已核验 \\
Baseline/对照 & 以论文或规范原文为准 & 不跨材料外推 \\
指标/对象 & 只使用当前来源可证明的信息 & E3 / 草案/未批准规范 \\
使用边界 & 可进入本方向演进图 & 不能替代更强证据 \\
\bottomrule
\end{tabularx}
}
    \end{minipage}
    \vspace{1mm}
  \end{minipage}}
  \end{column}
\end{columns}
\SourceFootnote{来源：ACAI: CCA accelerator execution｜证据等级：E3 / 草案/未批准规范}
\end{frame}


\begin{frame}[t,shrink=6]{ACAI: CCA accelerator execution：文章评价}
\DeckKicker{文章评价}
\SlideMeta{证据等级}{E3 / 草案/未批准规范}{来源状态：本地 PDF 已核验}
{\large\bfseries\color{Ink} 把 Arm CCA 讨论从 CPU/VM 推到 accelerator workflowow，是本方向的入口设计点\par}\vspace{1.2mm}
\begin{columns}[T,totalwidth=\textwidth]
  \begin{column}{0.45\textwidth}
    \fcolorbox{Rule}{Paper}{%
  \begin{minipage}[t]{0.97\linewidth}
    \vspace{1.1mm}
    \hspace{1.4mm}\begin{minipage}{0.92\linewidth}
      \PanelTitle{内容说明}
      \scriptsize \begin{itemize}
  \item 设计优势：把 Arm CCA 讨论从 CPU/VM 推到 accelerator workflowow，是本方向的入口设计点。
  \item 局限性：预印本不能代表标准化接口或量产设备能力，且驱动/runtime TCB 仍是主要工程风险。
  \item 商业化潜力：可服务 AI/HPC confidential offload；风险在设备厂商支持、驱动隔离和 attestation 生态不统一。
  \item 证据边界：E3 / Draft-not-ratified；本地 PDF 已核验。
\end{itemize}
    \end{minipage}
    \vspace{1mm}
  \end{minipage}}
  \end{column}
  \begin{column}{0.49\textwidth}
    \fcolorbox{Rule}{Panel}{%
  \begin{minipage}[t]{0.97\linewidth}
    \vspace{1.1mm}
    \hspace{1.4mm}\begin{minipage}{0.92\linewidth}
      \PanelTitle{文章评价}
      \scriptsize {\tiny
\begin{tabularx}{\linewidth}{@{}YY@{}}
\toprule
\textbf{维度} & \textbf{判断} \\
\midrule
设计优势 & 把 Arm CCA 讨论从 CPU/VM 推到 accelerator workflowow，是本方向的入口设计点。 \\
主要局限 & 预印本不能代表标准化接口或量产设备能力，且驱动/runtime TCB 仍是主要工程风险。 \\
商业落地 & 可服务 AI/HPC confidential offload；风险在设备厂商支持、驱动隔离和 attestation 生态不统一。 \\
讲稿定位 & 开创/基础入口; 证据强度 E3 \\
\bottomrule
\end{tabularx}
}
    \end{minipage}
    \vspace{1mm}
  \end{minipage}}
  \end{column}
\end{columns}
\SourceFootnote{来源：ACAI: CCA accelerator execution｜证据等级：E3 / 草案/未批准规范}
\end{frame}


\begin{frame}[t,shrink=6]{Devlore: device interrupt protection：内容摘要}
\DeckKicker{内容摘要}
\SlideMeta{证据等级}{E3 / 草案/未批准规范}{来源状态：本地 PDF 已核验}
{\large\bfseries\color{Ink} 关注 confidential VM 的 device interrupt protection，把 interrupt delivery 纳入机。\par}\vspace{1.2mm}
\begin{columns}[T,totalwidth=\textwidth]
  \begin{column}{0.45\textwidth}
    \fcolorbox{Rule}{Paper}{%
  \begin{minipage}[t]{0.97\linewidth}
    \vspace{1.1mm}
    \hspace{1.4mm}\begin{minipage}{0.92\linewidth}
      \PanelTitle{内容说明}
      \scriptsize \begin{itemize}
  \item 关注 confidential VM 的 device interrupt protection，把 interrupt delivery 纳入机密计算边界
  \item Devlore is a recent draft on device interrupt protection for confidential VMs and...
\end{itemize}
    \end{minipage}
    \vspace{1mm}
  \end{minipage}}
  \end{column}
  \begin{column}{0.49\textwidth}
    \fcolorbox{Rule}{Panel}{%
  \begin{minipage}[t]{0.97\linewidth}
    \vspace{1.1mm}
    \hspace{1.4mm}\begin{minipage}{0.92\linewidth}
      \PanelTitle{证据定位}
      \scriptsize {\tiny
\begin{tabularx}{\linewidth}{@{}YY@{}}
\toprule
\textbf{字段} & \textbf{内容} \\
\midrule
论文定位 & 代表性改进 \\
材料类型 & 系统论文 \\
证据等级 & E3 / 草案/未批准规范 \\
来源状态 & 本地 PDF 已核验 \\
\bottomrule
\end{tabularx}
}
    \end{minipage}
    \vspace{1mm}
  \end{minipage}}
  \end{column}
\end{columns}
\SourceFootnote{来源：Devlore: device interrupt protection｜证据等级：E3 / 草案/未批准规范}
\end{frame}


\begin{frame}[t,shrink=6]{Devlore: device interrupt protection：研究背景}
\DeckKicker{研究背景}
\SlideMeta{证据等级}{E3 / 草案/未批准规范}{来源状态：本地 PDF 已核验}
{\large\bfseries\color{Ink} 即便 DMA 受控，设备 interrupt 仍可能泄露 timing/ownership 信息，或被 host 用来破坏 confidential VM 的状态机\par}\vspace{1.2mm}
\begin{columns}[T,totalwidth=\textwidth]
  \begin{column}{0.45\textwidth}
    \fcolorbox{Rule}{Paper}{%
  \begin{minipage}[t]{0.97\linewidth}
    \vspace{1.1mm}
    \hspace{1.4mm}\begin{minipage}{0.92\linewidth}
      \PanelTitle{内容说明}
      \scriptsize \begin{itemize}
  \item 即便 DMA 受控，设备 interrupt 仍可能泄露 timing/ownership 信息，或被 host 用来破坏 confidential VM 的状态机
\end{itemize}
    \end{minipage}
    \vspace{1mm}
  \end{minipage}}
  \end{column}
  \begin{column}{0.49\textwidth}
    \fcolorbox{Rule}{Panel}{%
  \begin{minipage}[t]{0.97\linewidth}
    \vspace{1.1mm}
    \hspace{1.4mm}\begin{minipage}{0.92\linewidth}
      \PanelTitle{问题边界}
      \scriptsize {\tiny
\begin{tabularx}{\linewidth}{@{}YY@{}}
\toprule
\textbf{维度} & \textbf{说明} \\
\midrule
保护对象 & Arm CCA I/O、DMA、accelerator、interrupt \\
传统瓶颈 & 把设备访问、DMA、interrupt ownership 与 accelerator workflowow 纳入 CCA 边界。 \\
本文入口 & Devlore: device interrupt protection \\
证据限制 & E3 / 草案/未批准规范 \\
\bottomrule
\end{tabularx}
}
    \end{minipage}
    \vspace{1mm}
  \end{minipage}}
  \end{column}
\end{columns}
\SourceFootnote{来源：Devlore: device interrupt protection｜证据等级：E3 / 草案/未批准规范}
\end{frame}


\begin{frame}[t,shrink=6]{Devlore: device interrupt protection：解决方案}
\DeckKicker{解决方案}
\SlideMeta{证据等级}{E3 / 草案/未批准规范}{来源状态：本地 PDF 已核验}
{\large\bfseries\color{Ink} 将 interrupt delivery、device ownership 与 confidential boundary 绑定，使设备事件不再只。\par}\vspace{1.2mm}
\begin{columns}[T,totalwidth=\textwidth]
  \begin{column}{0.45\textwidth}
    \fcolorbox{Rule}{Paper}{%
  \begin{minipage}[t]{0.97\linewidth}
    \vspace{1.1mm}
    \hspace{1.4mm}\begin{minipage}{0.92\linewidth}
      \PanelTitle{内容说明}
      \scriptsize \begin{itemize}
  \item 将 interrupt delivery、device ownership 与 confidential boundary 绑定，使设备事件不再只是外围控制面问题
\end{itemize}
    \end{minipage}
    \vspace{1mm}
  \end{minipage}}
  \end{column}
  \begin{column}{0.49\textwidth}
    \fcolorbox{Rule}{Panel}{%
  \begin{minipage}[t]{0.97\linewidth}
    \vspace{1.1mm}
    \hspace{1.4mm}\begin{minipage}{0.92\linewidth}
      \PanelTitle{核心设计拆解}
      \scriptsize \begin{itemize}
  \item 01. Confidential-device interrupt ownership
  \item 02. Interrupt delivery mediation
  \item 03. VM/device lifecycle coupling
\end{itemize}
    \end{minipage}
    \vspace{1mm}
  \end{minipage}}
  \end{column}
\end{columns}
\SourceFootnote{来源：Devlore: device interrupt protection｜证据等级：E3 / 草案/未批准规范}
\end{frame}


\begin{frame}[t,shrink=6]{Devlore: device interrupt protection：实验结果}
\DeckKicker{实验结果}
\SlideMeta{证据等级}{E3 / 草案/未批准规范}{来源状态：本地 PDF 已核验}
{\large\bfseries\color{Ink} 本地已核验 arXiv PDF 提供 device-interrupt protection 设计/原型证据\par}\vspace{1.2mm}
\begin{columns}[T,totalwidth=\textwidth]
  \begin{column}{0.45\textwidth}
    \fcolorbox{Rule}{Paper}{%
  \begin{minipage}[t]{0.97\linewidth}
    \vspace{1.1mm}
    \hspace{1.4mm}\begin{minipage}{0.92\linewidth}
      \PanelTitle{内容说明}
      \scriptsize \begin{itemize}
  \item 本地已核验 arXiv PDF 提供 device-interrupt protection 设计/原型证据
  \item 本报告 只写机制方向和证据状态，不写未复核数字或成熟部署结论
\end{itemize}
    \end{minipage}
    \vspace{1mm}
  \end{minipage}}
  \end{column}
  \begin{column}{0.49\textwidth}
    \fcolorbox{Rule}{Panel}{%
  \begin{minipage}[t]{0.97\linewidth}
    \vspace{1.1mm}
    \hspace{1.4mm}\begin{minipage}{0.92\linewidth}
      \PanelTitle{实验/证据边界}
      \scriptsize {\tiny
\begin{tabularx}{\linewidth}{@{}YYY@{}}
\toprule
\textbf{对象} & \textbf{可支撑} & \textbf{边界} \\
\midrule
数据/来源 & Devlore: device interrupt protection & 本地 PDF 已核验 \\
Baseline/对照 & 以论文或规范原文为准 & 不跨材料外推 \\
指标/对象 & 只使用当前来源可证明的信息 & E3 / 草案/未批准规范 \\
使用边界 & 可进入本方向演进图 & 不能替代更强证据 \\
\bottomrule
\end{tabularx}
}
    \end{minipage}
    \vspace{1mm}
  \end{minipage}}
  \end{column}
\end{columns}
\SourceFootnote{来源：Devlore: device interrupt protection｜证据等级：E3 / 草案/未批准规范}
\end{frame}


\begin{frame}[t,shrink=6]{Devlore: device interrupt protection：文章评价}
\DeckKicker{文章评价}
\SlideMeta{证据等级}{E3 / 草案/未批准规范}{来源状态：本地 PDF 已核验}
{\large\bfseries\color{Ink} 明确指出 interrupt 是 CCA device path 的核心安全面，不应只讨论 DMA\par}\vspace{1.2mm}
\begin{columns}[T,totalwidth=\textwidth]
  \begin{column}{0.45\textwidth}
    \fcolorbox{Rule}{Paper}{%
  \begin{minipage}[t]{0.97\linewidth}
    \vspace{1.1mm}
    \hspace{1.4mm}\begin{minipage}{0.92\linewidth}
      \PanelTitle{内容说明}
      \scriptsize \begin{itemize}
  \item 设计优势：明确指出 interrupt 是 CCA device path 的核心安全面，不应只讨论 DMA。
  \item 局限性：Draft evidence，落地需硬件、firmware、OS、hypervisor 和 interrupt controller 协同。
  \item 商业化潜力：对云端直通设备和加速器租户隔离很关键；风险在 legacy 设备兼容和中断虚拟化复杂性。
  \item 证据边界：E3 / Draft-not-ratified；本地 PDF 已核验。
\end{itemize}
    \end{minipage}
    \vspace{1mm}
  \end{minipage}}
  \end{column}
  \begin{column}{0.49\textwidth}
    \fcolorbox{Rule}{Panel}{%
  \begin{minipage}[t]{0.97\linewidth}
    \vspace{1.1mm}
    \hspace{1.4mm}\begin{minipage}{0.92\linewidth}
      \PanelTitle{文章评价}
      \scriptsize {\tiny
\begin{tabularx}{\linewidth}{@{}YY@{}}
\toprule
\textbf{维度} & \textbf{判断} \\
\midrule
设计优势 & 明确指出 interrupt 是 CCA device path 的核心安全面，不应只讨论 DMA。 \\
主要局限 & Draft evidence，落地需硬件、firmware、OS、hypervisor 和 interrupt controller 协同。 \\
商业落地 & 对云端直通设备和加速器租户隔离很关键；风险在 legacy 设备兼容和中断虚拟化复杂性。 \\
讲稿定位 & 代表性改进; 证据强度 E3 \\
\bottomrule
\end{tabularx}
}
    \end{minipage}
    \vspace{1mm}
  \end{minipage}}
  \end{column}
\end{columns}
\SourceFootnote{来源：Devlore: device interrupt protection｜证据等级：E3 / 草案/未批准规范}
\end{frame}


\begin{frame}[t,shrink=6]{Accelerator TEE SoK for CCA I/O：内容摘要}
\DeckKicker{内容摘要}
\SlideMeta{证据等级}{E2 / Background substrate}{来源状态：本地 PDF 已核验}
{\large\bfseries\color{Ink} 为 accelerator/device TEE 提供 taxonomy，帮助组织 GPU/NPU/DPU/SmartNIC 的身份、内存路径、runtime 与 attestation 问题\par}\vspace{1.2mm}
\begin{columns}[T,totalwidth=\textwidth]
  \begin{column}{0.45\textwidth}
    \fcolorbox{Rule}{Paper}{%
  \begin{minipage}[t]{0.97\linewidth}
    \vspace{1.1mm}
    \hspace{1.4mm}\begin{minipage}{0.92\linewidth}
      \PanelTitle{内容说明}
      \scriptsize \begin{itemize}
  \item 为 accelerator/device TEE 提供 taxonomy，帮助组织 GPU/NPU/DPU/SmartNIC 的身份、内存路径、runtime 与。
  \item Accelerator TEE SoK provides current taxonomy for device identity, memory paths, runtime...
\end{itemize}
    \end{minipage}
    \vspace{1mm}
  \end{minipage}}
  \end{column}
  \begin{column}{0.49\textwidth}
    \fcolorbox{Rule}{Panel}{%
  \begin{minipage}[t]{0.97\linewidth}
    \vspace{1.1mm}
    \hspace{1.4mm}\begin{minipage}{0.92\linewidth}
      \PanelTitle{证据定位}
      \scriptsize {\tiny
\begin{tabularx}{\linewidth}{@{}YY@{}}
\toprule
\textbf{字段} & \textbf{内容} \\
\midrule
论文定位 & 当前 SOTA/补充边界 \\
材料类型 & SoK/综述 \\
证据等级 & E2 / Background substrate \\
来源状态 & 本地 PDF 已核验 \\
\bottomrule
\end{tabularx}
}
    \end{minipage}
    \vspace{1mm}
  \end{minipage}}
  \end{column}
\end{columns}
\SourceFootnote{来源：Accelerator TEE SoK for CCA I/O｜证据等级：E2 / Background substrate}
\end{frame}


\begin{frame}[t,shrink=6]{Accelerator TEE SoK for CCA I/O：研究背景}
\DeckKicker{研究背景}
\SlideMeta{证据等级}{E2 / Background substrate}{来源状态：本地 PDF 已核验}
{\large\bfseries\color{Ink} 设备 TEE 机制比 CPU TEE 更分散\par}\vspace{1.2mm}
\begin{columns}[T,totalwidth=\textwidth]
  \begin{column}{0.45\textwidth}
    \fcolorbox{Rule}{Paper}{%
  \begin{minipage}[t]{0.97\linewidth}
    \vspace{1.1mm}
    \hspace{1.4mm}\begin{minipage}{0.92\linewidth}
      \PanelTitle{内容说明}
      \scriptsize \begin{itemize}
  \item 设备 TEE 机制比 CPU TEE 更分散
  \item 不同 accelerator 的队列、固件、驱动和本地 root 差异很大
\end{itemize}
    \end{minipage}
    \vspace{1mm}
  \end{minipage}}
  \end{column}
  \begin{column}{0.49\textwidth}
    \fcolorbox{Rule}{Panel}{%
  \begin{minipage}[t]{0.97\linewidth}
    \vspace{1.1mm}
    \hspace{1.4mm}\begin{minipage}{0.92\linewidth}
      \PanelTitle{问题边界}
      \scriptsize {\tiny
\begin{tabularx}{\linewidth}{@{}YY@{}}
\toprule
\textbf{维度} & \textbf{说明} \\
\midrule
保护对象 & Arm CCA I/O、DMA、accelerator、interrupt \\
传统瓶颈 & 把设备访问、DMA、interrupt ownership 与 accelerator workflowow 纳入 CCA 边界。 \\
本文入口 & Accelerator TEE SoK for CCA I/O \\
证据限制 & E2 / Background substrate \\
\bottomrule
\end{tabularx}
}
    \end{minipage}
    \vspace{1mm}
  \end{minipage}}
  \end{column}
\end{columns}
\SourceFootnote{来源：Accelerator TEE SoK for CCA I/O｜证据等级：E2 / Background substrate}
\end{frame}


\begin{frame}[t,shrink=6]{Accelerator TEE SoK for CCA I/O：解决方案}
\DeckKicker{解决方案}
\SlideMeta{证据等级}{E2 / Background substrate}{来源状态：本地 PDF 已核验}
{\large\bfseries\color{Ink} 用 taxonomy 区分设备身份、内存路径、队列、runtime、driver TCB 与 attestation，作为 CCA I/O 方向的。\par}\vspace{1.2mm}
\begin{columns}[T,totalwidth=\textwidth]
  \begin{column}{0.45\textwidth}
    \fcolorbox{Rule}{Paper}{%
  \begin{minipage}[t]{0.97\linewidth}
    \vspace{1.1mm}
    \hspace{1.4mm}\begin{minipage}{0.92\linewidth}
      \PanelTitle{内容说明}
      \scriptsize \begin{itemize}
  \item 用 taxonomy 区分设备身份、内存路径、队列、runtime、driver TCB 与 attestation，作为 CCA I/O 方向的背景基底
\end{itemize}
    \end{minipage}
    \vspace{1mm}
  \end{minipage}}
  \end{column}
  \begin{column}{0.49\textwidth}
    \fcolorbox{Rule}{Panel}{%
  \begin{minipage}[t]{0.97\linewidth}
    \vspace{1.1mm}
    \hspace{1.4mm}\begin{minipage}{0.92\linewidth}
      \PanelTitle{核心设计拆解}
      \scriptsize \begin{itemize}
  \item 01. Accelerator/device TEE design taxonomy
  \item 02. Device identity and attestation vocabulary
  \item 03. Memory, queue, runtime and driver TCB dimensions
\end{itemize}
    \end{minipage}
    \vspace{1mm}
  \end{minipage}}
  \end{column}
\end{columns}
\SourceFootnote{来源：Accelerator TEE SoK for CCA I/O｜证据等级：E2 / Background substrate}
\end{frame}


\begin{frame}[t,shrink=6]{Accelerator TEE SoK for CCA I/O：实验结果}
\DeckKicker{实验结果}
\SlideMeta{证据等级}{E2 / Background substrate}{来源状态：本地 PDF 已核验}
{\large\bfseries\color{Ink} 规范/Survey，无新实验；Accelerator TEE SoK for CCA I/O 只支撑术语、taxonomy、接口语义或证据边界。\par}\vspace{1.2mm}
\begin{columns}[T,totalwidth=\textwidth]
  \begin{column}{0.45\textwidth}
    \fcolorbox{Rule}{Paper}{%
  \begin{minipage}[t]{0.97\linewidth}
    \vspace{1.1mm}
    \hspace{1.4mm}\begin{minipage}{0.92\linewidth}
      \PanelTitle{内容说明}
      \scriptsize \begin{itemize}
  \item 本页不展示独立系统实验，也不把二手归纳写成一手结果。
  \item 可支撑：Accelerator TEE SoK provides current taxonomy for device identity, memory paths...
  \item 证据等级：E2 / Background substrate；来源状态：本地 PDF 已核验。
  \item 涉及具体平台、协议状态或数值时，转到原始论文、官方规范或厂商材料核验。
\end{itemize}
    \end{minipage}
    \vspace{1mm}
  \end{minipage}}
  \end{column}
  \begin{column}{0.49\textwidth}
    \fcolorbox{Rule}{Panel}{%
  \begin{minipage}[t]{0.97\linewidth}
    \vspace{1.1mm}
    \hspace{1.4mm}\begin{minipage}{0.92\linewidth}
      \PanelTitle{实验/证据边界}
      \scriptsize {\tiny
\begin{tabularx}{\linewidth}{@{}YYY@{}}
\toprule
\textbf{对象} & \textbf{可支撑} & \textbf{边界} \\
\midrule
数据/来源 & Accelerator TEE SoK for CCA I/O & 本地 PDF 已核验 \\
Baseline/对照 & 以论文或规范原文为准 & 不跨材料外推 \\
指标/对象 & 只使用当前来源可证明的信息 & E2 / Background substrate \\
使用边界 & 可进入本方向演进图 & 不能替代更强证据 \\
\bottomrule
\end{tabularx}
}
    \end{minipage}
    \vspace{1mm}
  \end{minipage}}
  \end{column}
\end{columns}
\SourceFootnote{来源：Accelerator TEE SoK for CCA I/O｜证据等级：E2 / Background substrate}
\end{frame}


\begin{frame}[t,shrink=6]{Accelerator TEE SoK for CCA I/O：文章评价}
\DeckKicker{文章评价}
\SlideMeta{证据等级}{E2 / Background substrate}{来源状态：本地 PDF 已核验}
{\large\bfseries\color{Ink} 评价：Accelerator TEE SoK for CCA I/O 适合做 taxonomy/规范入口；规范/Survey，无新实验，不能替代一手系统证据。\par}\vspace{1.2mm}
\begin{columns}[T,totalwidth=\textwidth]
  \begin{column}{0.45\textwidth}
    \fcolorbox{Rule}{Paper}{%
  \begin{minipage}[t]{0.97\linewidth}
    \vspace{1.1mm}
    \hspace{1.4mm}\begin{minipage}{0.92\linewidth}
      \PanelTitle{内容说明}
      \scriptsize \begin{itemize}
  \item 设计优势：能把 accelerator TEE 设计空间系统化，避免只围绕 CPU Realm 叙事。
  \item 局限性：不是 Arm CCA 专用机制证据，也不是某个设备的生产部署证明。
  \item 商业化潜力：有助于产品需求拆解和 vendor 对标；风险在 taxonomy 与具体硬件能力之间仍需二次验证。
  \item 规范/Survey，无新实验；具体平台结论转到一手材料核验。
\end{itemize}
    \end{minipage}
    \vspace{1mm}
  \end{minipage}}
  \end{column}
  \begin{column}{0.49\textwidth}
    \fcolorbox{Rule}{Panel}{%
  \begin{minipage}[t]{0.97\linewidth}
    \vspace{1.1mm}
    \hspace{1.4mm}\begin{minipage}{0.92\linewidth}
      \PanelTitle{文章评价}
      \scriptsize {\tiny
\begin{tabularx}{\linewidth}{@{}YY@{}}
\toprule
\textbf{维度} & \textbf{判断} \\
\midrule
设计优势 & 能把 accelerator TEE 设计空间系统化，避免只围绕 CPU Realm 叙事。 \\
主要局限 & 不是 Arm CCA 专用机制证据，也不是某个设备的生产部署证明。 \\
商业落地 & 有助于产品需求拆解和 vendor 对标；风险在 taxonomy 与具体硬件能力之间仍需二次验证。 \\
讲稿定位 & 当前 SOTA/补充边界; 证据强度 E2 \\
\bottomrule
\end{tabularx}
}
    \end{minipage}
    \vspace{1mm}
  \end{minipage}}
  \end{column}
\end{columns}
\SourceFootnote{来源：Accelerator TEE SoK for CCA I/O｜证据等级：E2 / Background substrate}
\end{frame}


\begin{frame}[t,shrink=6]{Arm CCA I/O、DMA、accelerator、interrupt：方向总结}
\DeckKicker{方向总结}
\SlideMeta{证据等级}{方向综述}{来源状态：公开材料综合}
{\large\bfseries\color{Ink} Arm CCA I/O、DMA、accelerator、interrupt 的结论是：三篇主讲材料共同给出技术演进，但仍要用证据等级约束每个结论。\par}\vspace{1.2mm}
\begin{columns}[T,totalwidth=\textwidth]
  \begin{column}{0.45\textwidth}
    \fcolorbox{Rule}{Paper}{%
  \begin{minipage}[t]{0.97\linewidth}
    \vspace{1.1mm}
    \hspace{1.4mm}\begin{minipage}{0.92\linewidth}
      \PanelTitle{内容说明}
      \scriptsize \begin{itemize}
  \item ACAI: CCA accelerator execution：开创/基础入口，E3 / 草案/未批准规范。
  \item Devlore: device interrupt protection：代表性改进，E3 / 草案/未批准规范。
  \item Accelerator TEE SoK for CCA I/O：当前 SOTA/补充边界，E2 / Background substrate。
  \item 本方向结论：先用三篇主讲材料建立演进关系，再用证据等级控制机制、实验和商业化结论。
\end{itemize}
    \end{minipage}
    \vspace{1mm}
  \end{minipage}}
  \end{column}
  \begin{column}{0.49\textwidth}
    \fcolorbox{Rule}{Panel}{%
  \begin{minipage}[t]{0.97\linewidth}
    \vspace{1.1mm}
    \hspace{1.4mm}\begin{minipage}{0.92\linewidth}
      \PanelTitle{本方向方法对比}
      \scriptsize {\tiny
\begin{tabularx}{\linewidth}{@{}YYY@{}}
\toprule
\textbf{论文} & \textbf{定位} & \textbf{选择理由} \\
\midrule
ACAI: CCA accelerator execution & 开创/基础入口 & ACAI is the early design point for extending Arm CCA toward accelerator... \\
Devlore: device interrupt protection & 代表性改进 & Devlore is a recent draft on device interrupt protection for... \\
Accelerator TEE SoK for CCA I/O & 当前 SOTA/补充边界 & Accelerator TEE SoK provides current taxonomy for device identity... \\
\bottomrule
\end{tabularx}
}
    \end{minipage}
    \vspace{1mm}
  \end{minipage}}
  \end{column}
\end{columns}
\SourceFootnote{来源：论文原文、官方规范或公开材料｜证据等级：见本页 badge}
\end{frame}
